\documentclass[10pt]{article}
\usepackage{fontspec}
\setmainfont{DejaVu Sans}
\usepackage[paperwidth=17in,paperheight=11in,margin=0in]{geometry}
\usepackage{tikz}
\usetikzlibrary{arrows.meta}
\pagestyle{empty}
\setlength{\parindent}{0pt}
\setlength{\topskip}{0pt}

\tikzset{
  wire/.style={draw,line width=.75pt,line cap=round},
  sym/.style={draw,line width=.85pt,line cap=round,line join=round},
  mech/.style={draw,line width=.55pt,densely dashed},
  frame/.style={draw,line width=.6pt},
  divider/.style={draw=black!55,line width=.45pt},
  title/.style={font=\fontsize{18}{21}\selectfont\bfseries},
  sheettitle/.style={font=\fontsize{14}{17}\selectfont},
  head/.style={font=\fontsize{10}{12}\selectfont\bfseries},
  label/.style={font=\fontsize{9}{11}\selectfont},
  pin/.style={font=\fontsize{8}{10}\selectfont},
  wireid/.style={font=\fontsize{8}{9}\selectfont,inner sep=1pt},
  note/.style={font=\fontsize{9}{13}\selectfont,align=left,text width=101mm,
    inner sep=0pt,anchor=north west}
}

\newcommand{\sheetframe}[2]{%
  \path[use as bounding box] (0,0) rectangle (169.9,109.9);
  \draw[frame] (4,4) rectangle (166,106);
  \draw[frame] (4,93)--(166,93) (4,12)--(166,12);
  \draw[divider] (116,16)--(116,90);
  \node[title,anchor=west] at (7,102) {SƠ ĐỒ ĐIỆN ĐỘNG CƠ 3 PHA};
  \node[sheettitle,anchor=west] at (7,97.4) {#2};
  \node[head,anchor=east] at (163,102) {BẢN NHÁP V0.4};
  \node[label,anchor=east] at (163,97.4) {KHÔNG DÙNG THI CÔNG};
  \node[label,anchor=west] at (7,8.7) {M3P-#1 \quad #2};
  \node[pin,anchor=west] at (7,5.8) {27-11-2026 \quad\textbar\quad 17 × 11 inch};
  \node[head,anchor=east] at (163,8) {TỜ #1 / 002};
  \node[head,anchor=west] at (119,88.4) {GHI CHÚ THIẾT KẾ};
  \draw[divider] (120,85.5)--(162,85.5);
}
\newcommand{\notegroup}[3]{%
  \node[head,anchor=west] at (119,#1) {#2};
  \node[note] at (120,{#1-2.2}) {#3};
}
\newcommand{\terminal}[2]{\draw[sym,fill=white] (#1,#2) circle[radius=.42];}
\newcommand{\junction}[2]{\fill (#1,#2) circle[radius=.32];}
\newcommand{\earth}[2]{%
  \begin{scope}[shift={(#1,#2)}]
    \draw[sym] (0,0)--(0,-1.4) (-1.6,-1.4)--(1.6,-1.4)
      (-1.05,-2)--(1.05,-2) (-.5,-2.6)--(.5,-2.6);
  \end{scope}
}
% Each contact has terminals at x = -3 and +3.
\newcommand{\opencontact}{%
  \draw[sym] (-2.65,.16)--(2.2,2);
  \terminal{-3}{0}\terminal{3}{0}
}
\newcommand{\closedcontact}{%
  \draw[sym] (-2.65,0)--(2.65,0) (1.1,-.9)--(2.55,.9);
  \terminal{-3}{0}\terminal{3}{0}
}
\newcommand{\pushactuator}{%
  \draw[mech] (0,2.95)--(0,1.25);
  \draw[sym] (0,2.95)--(0,3.8) (-1.55,3.8)--(1.55,3.8);
}
\newcommand{\emergencyactuator}{%
  \draw[mech] (0,2.95)--(0,.35);
  \draw[sym] (0,2.95)--(0,3.7) (-1.8,3.7)--(1.8,3.7)
    (-1.8,3.7) arc[start angle=180,end angle=0,x radius=1.8,y radius=.95];
}
\newcommand{\lamp}{%
  \draw[sym] (0,0) circle[radius=1.7]
    (-1.2,-1.2)--(1.2,1.2) (-1.2,1.2)--(1.2,-1.2);
}
\newcommand{\wirelabel}[3]{\node[wireid,anchor=south] at (#1,#2) {#3};}
\newcommand{\netflag}[3]{%
  \draw[sym,fill=white] (#1-.9,#2+.95)--(#1,#2)--(#1+.9,#2+.95)--cycle;
  \node[pin,anchor=south] at (#1,#2+1.1) {#3};
}

\begin{document}
\noindent
\begin{tikzpicture}[x=.1in,y=.1in]
\sheetframe{001}{MẠCH ĐỘNG LỰC 3 PHA}
\node[head,anchor=west] at (8,88.4) {KHỞI ĐỘNG TRỰC TIẾP DOL};
\node[label,anchor=west] at (8,84.7) {Nguồn 3 pha và PE \quad\textbar\quad Ký hiệu ở trạng thái chưa tác động};

\node[head] at (31,79) {QF1};
\node[pin] at (31,76.5) {CB 3P};
\node[head] at (53,79) {KM1};
\node[pin] at (53,76.5) {Contactor 3P};
\node[head] at (73,79) {FR1};
\node[pin] at (73,76.5) {Rơ-le nhiệt 3P};
\node[head] at (103,79) {M1};

% Phase changes are explicit, with common mechanical actuation per device.
\foreach \x in {31,53} {
  \draw[mech] (\x,60.15)--(\x,72.15) (\x,48.15)--(\x,60.15);
}
\draw[mech] (73,61.65)--(73,71.6) (73,49.65)--(73,59.6);
\foreach \y/\ph/\inwire/\cbwire/\kmwire/\outwire/\linepin/\loadpin/\motorpin in
  {71/L1/101/111/121/131/1L1/2T1/U1,
   59/L2/102/112/122/132/3L2/4T2/V1,
   47/L3/103/113/123/133/5L3/6T3/W1} {
  \node[label,anchor=east] at (13,\y) {\ph};
  \draw[wire] (15,\y)--(28,\y) (34,\y)--(50,\y)
    (56,\y)--(70,\y) (76,\y)--(91,\y);
  \terminal{15}{\y}\terminal{91}{\y}
  \begin{scope}[shift={(31,\y)}]
    \draw[sym] (-2.65,.16)--(2.2,2)
      (2.35,-.65)--(3.65,.65) (2.35,.65)--(3.65,-.65);
    \terminal{-3}{0}
  \end{scope}
  \begin{scope}[shift={(53,\y)}]\opencontact\end{scope}
  \begin{scope}[shift={(73,\y)}]
    \draw[sym] (-3,0)--(-2.2,0)--(-2.2,1.65)--(2.2,1.65)--(2.2,0)--(3,0);
  \end{scope}
  \wirelabel{21}{\y+.75}{\inwire}
  \wirelabel{42}{\y+.75}{\cbwire}
  \wirelabel{63}{\y+.75}{\kmwire}
  \wirelabel{83}{\y+.75}{\outwire}
  \node[pin,anchor=north] at (28,\y-.75) {\linepin};
  \node[pin,anchor=north] at (34,\y-.75) {\loadpin};
  \node[pin,anchor=north] at (50,\y-.75) {\linepin};
  \node[pin,anchor=north] at (56,\y-.75) {\loadpin};
  \node[pin,anchor=east] at (90,\y-2.3) {\motorpin};
}
\draw[wire] (91.42,71)--(96,71)--(96,63.2426)
  (91.42,59)--(97,59)
  (91.42,47)--(96,47)--(96,54.7574);
\draw[sym] (103,59) circle[radius=6];
% The outer phase wires end on the motor symbol.
\draw[wire] (96,63.2426)--(98.7574,63.2426)
  (96,54.7574)--(98.7574,54.7574);
\node[head] at (103,60.6) {M};
\node[label] at (103,57.2) {3\textasciitilde{}};
\node[label,anchor=east] at (13,34) {PE};
\draw[wire] (15,34)--(103,34)--(103,53);
\terminal{15}{34}
\wirelabel{21}{34.75}{104}
\junction{103}{34}\earth{103}{34}
\node[pin,anchor=east] at (101,38) {Vỏ động cơ};

\draw[divider] (9,26)--(112,26);
\node[label,anchor=west] at (8,22) {Tờ 002 dùng cuộn KM1 (A1, A2) và tiếp điểm phụ FR1 (95, 96).};
\node[pin,anchor=west] at (8,18) {Tham chiếu ký hiệu IEC và sơ đồ DOL của ABB, tài liệu 1SFC133001C0201.};

\notegroup{81.5}{01 \quad BẢO VỆ NGUỒN}{%
  Chọn QF1 theo dòng làm việc, khả năng cắt và điều kiện phối hợp bảo vệ của tủ.}
\notegroup{65}{02 \quad BẢO VỆ QUÁ TẢI}{%
  FR1 đặt sau KM1. Chỉnh dòng theo nhãn động cơ và hướng dẫn của nhà sản xuất.}
\notegroup{48.5}{03 \quad ĐẤU NỐI ĐỘNG CƠ}{%
  Ba pha vào U1, V1 và W1. Dây PE số 104 đi trực tiếp đến vỏ động cơ.
  Xác định cách đấu sao hoặc tam giác theo nhãn M1 và điện áp nguồn.}
\notegroup{32}{04 \quad ĐÁNH SỐ DÂY}{%
  101–103 trước QF1, 111–113 sau QF1, 121–123 sau KM1 và 131–133 đến M1.}
\end{tikzpicture}
\par\newpage
\noindent
\begin{tikzpicture}[x=.1in,y=.1in]
\sheetframe{002}{MẠCH ĐIỀU KHIỂN 24 VDC VÀ NÚT NHẤN}
\node[head,anchor=west] at (8,88.4) {NGUỒN ĐIỀU KHIỂN VÀ CÁC NGÕ VÀO RA};
\node[pin,anchor=west] at (8,83.7) {230 VAC phụ trợ};

% AC/DC converter and protected distribution, with repeated net references.
\draw[wire] (15,79.5)--(29,79.5) (15,74.5)--(29,74.5);
\node[label,anchor=east] at (14.5,79.5) {L};
\node[label,anchor=east] at (14.5,74.5) {N};
\draw[sym] (29,72) rectangle (43,82) (29,72)--(43,82);
\node[head,anchor=south] at (36,82.1) {PS1};
\node[label] at (32.5,78.5) {\textasciitilde{}};
\draw[sym] (37.9,75.9)--(40.9,75.9);
\draw[sym] (37.9,75)--(38.6,75) (39.05,75)--(39.75,75) (40.2,75)--(40.9,75);
\draw[wire] (43,79.5)--(53,79.5) (59,79.5)--(77,79.5);
\draw[sym] (53,78.25) rectangle (59,80.75) (53,79.5)--(59,79.5);
\node[head,anchor=south] at (56,81.1) {F1};
\draw[wire] (43,74.5)--(68,74.5);
\node[pin,anchor=west] at (43,81.4) {+};
\node[pin,anchor=west] at (43,72.6) {−};
\wirelabel{48}{80.05}{201}
\wirelabel{65}{80.05}{202}
\wirelabel{55}{75.2}{200}
\draw[sym,fill=white] (77,78.6)--(77.95,79.5)--(77,80.4)--cycle
  (68,73.6)--(68.95,74.5)--(68,75.4)--cycle;
\node[pin,anchor=west] at (78.6,79.5) {+24 VDC};
\node[pin,anchor=west] at (69.6,74.5) {0 VDC};

% Left distribution rail is net 202. Each return has its own module pin.
\draw[wire] (15,68.7)--(15,18.7);
\netflag{15}{68.7}{202}
\foreach \y in {65,47,37.5,28,18.7} {\junction{15}{\y}}
\draw[sym] (88,15.5) rectangle (112,70.5);
\node[head] at (100,65.8) {A1 \quad DEVICE\_1M};
\node[pin,align=center] at (100,54.1) {DO NPN kéo xuống 0 V\\DI nhận mức +24 V};
\draw[wire] (95,70.5)--(95,74.2) (105,70.5)--(105,74.2);
\netflag{95}{74.2}{202}
\netflag{105}{74.2}{200}
\node[pin,anchor=north] at (95,70.3) {V+};
\node[pin,anchor=north] at (105,70.3) {0V/COM};

% Hard-wired stop chain. D1 has cathode at coil A1 (left).
\draw[wire] (15,65)--(25,65) (31,65)--(43,65) (49,65)--(61,65)
  (69,65)--(80,65)--(80,60)--(88,60);
\begin{scope}[shift={(28,65)}]\closedcontact\emergencyactuator\end{scope}
\begin{scope}[shift={(46,65)}]\closedcontact\end{scope}
\draw[sym] (61,63.4) rectangle (69,66.6);
\node[head,anchor=west] at (30.1,69.5) {S0};
\node[pin,anchor=north,align=center] at (28,63.9) {Dừng khẩn NC\\11–12*};
\node[head] at (46,68.8) {FR1};
\node[pin,anchor=north] at (46,63.9) {95–96 NC};
\node[head] at (65,69.5) {KM1};
\node[pin,anchor=north east] at (61,64.2) {A1};
\node[pin,anchor=north west] at (69,64.2) {A2};
\wirelabel{37}{65.7}{203}
\wirelabel{54}{65.7}{204}
\wirelabel{75.8}{65.7}{205}
\draw[wire] (57,65)--(57,56.5)--(63,56.5)
  (67.5,56.5)--(74,56.5)--(74,65);
\junction{57}{65}\junction{74}{65}
\draw[sym] (63,54.9)--(63,58.1)
  (67.5,55)--(63,56.5)--(67.5,58)--cycle;
\node[pin,anchor=north] at (65.5,54.5) {D1};

% Start and Stop inputs and indicator lamps with separate NPN outputs.
\draw[wire] (15,47)--(36,47) (42,47)--(88,47)
  (15,37.5)--(36,37.5) (42,37.5)--(88,37.5)
  (15,28)--(37.3,28) (40.7,28)--(88,28)
  (15,18.7)--(37.3,18.7) (40.7,18.7)--(88,18.7);
\begin{scope}[shift={(39,47)}]\opencontact\pushactuator\end{scope}
\begin{scope}[shift={(39,37.5)}]\closedcontact\pushactuator\end{scope}
\begin{scope}[shift={(39,28)}]\lamp\end{scope}
\begin{scope}[shift={(39,18.7)}]\lamp\end{scope}
\node[head,anchor=west] at (42.2,50.3) {S1 \enspace Khởi động};
\node[pin,anchor=north] at (39,45.8) {13–14* \quad NO};
\node[head,anchor=west] at (42.2,40.8) {S2 \enspace Dừng};
\node[pin,anchor=north] at (39,36.3) {11–12* \quad NC};
\node[pin,anchor=west] at (43.5,30.7) {H1 \quad LED nút S1};
\node[pin,anchor=west] at (43.5,21.4) {H2 \quad LED nút S2};
\foreach \y in {28,18.7} {
  \node[pin,anchor=east] at (37,\y+1.4) {+};
  \node[pin,anchor=west] at (41,\y+1.4) {−};
}
\foreach \y/\wid/\port in {
    47/213/DI\_START,37.5/216/DI\_STOP,
    28/222/DO\_LED\_START,18.7/232/DO\_LED\_STOP} {
  \wirelabel{72}{\y+.8}{\wid}
  \terminal{88}{\y}
  \node[pin,anchor=west] at (88.5,\y) {\texttt{\port}};
}
\terminal{88}{60}
\node[pin,anchor=west] at (88.5,60) {\texttt{DO\_KM1}};

\notegroup{81.5}{01 \quad DỪNG KHẨN}{%
  S0 (NC) và FR1 (95–96 NC) cắt trực tiếp nguồn cuộn KM1.
  Thẩm định chức năng dừng khẩn theo mức rủi ro máy.}
\notegroup{66}{02 \quad CUỘN KM1 VÀ DIODE D1}{%
  Cuộn KM1 dùng 24 VDC. Cathode D1 ở A1 cuộn KM1, anode ở A2.
  Ngõ DO\_KM1 phải chịu được dòng cuộn.}
\notegroup{50.5}{03 \quad NÚT NHẤN VÀ ĐÈN}{%
  S1 và S2 đưa lệnh vào phần mềm. H1 và H2 là cụm đèn LED 24 VDC.
  * Đối chiếu chân nút theo thiết bị chọn thực tế.}
\notegroup{35}{04 \quad NGUỒN VÀ KHỞI ĐỘNG LẠI}{%
  Xác nhận nguồn 230 VAC phụ trợ, sơ đồ chân DEVICE\_1M và điều kiện
  ngăn tự khởi động lại sau khi nhả S0 hoặc có điện trở lại.}
\end{tikzpicture}
\par
\end{document}
